\documentclass[11pt,a4paper]{article}
\usepackage[utf8]{inputenc}
\usepackage[T1]{fontenc}
\usepackage{lmodern}
\usepackage{amsmath}
\usepackage[margin=2cm]{geometry}
\usepackage{enumitem}
\DeclareUnicodeCharacter{2190}{\ensuremath{\leftarrow}}
\pagestyle{plain}
\setlength{\parindent}{0pt}
\setlength{\parskip}{0pt}
\begin{document}
{\Large\bfseries Tarea de lenguaje algorítmico}
\vspace{12pt}
\begin{enumerate}[label=\textbf{\arabic*.},leftmargin=*,itemsep=11pt,topsep=0pt]
\item \begin{minipage}[t]{\linewidth}
¿Qué tipo de dato es el más adecuado para guardar el promedio de notas de un estudiante?

\vspace{4pt}
\begin{tabular}{@{}p{0.47\linewidth}p{0.47\linewidth}@{}}
a)~Entero & b)~Cadena \\
c)~Real & d)~Booleano \\
\end{tabular}
\end{minipage}

\item \begin{minipage}[t]{\linewidth}
¿Qué afirmación sobre una constante es correcta?

\vspace{4pt}
\begin{tabular}{@{}p{0.47\linewidth}p{0.47\linewidth}@{}}
a)~Su valor cambia cada vez que se repite un bucle & b)~Se asigna una sola vez y no puede modificarse durante la ejecución \\
c)~Solo puede guardar texto & d)~Se escribe con camelCase por convención \\
\end{tabular}
\end{minipage}

\item \begin{minipage}[t]{\linewidth}
¿Cuál es el resultado de (6 + 2) * 3 - 4?

\vspace{4pt}
\begin{tabular}{@{}p{0.47\linewidth}p{0.47\linewidth}@{}}
a)~24 & b)~20 \\
c)~12 & d)~8 \\
\end{tabular}
\end{minipage}

\item \begin{minipage}[t]{\linewidth}
¿Cuál es el resultado de 2 \textasciicircum{} 3 * 2?

\vspace{4pt}
\begin{tabular}{@{}p{0.47\linewidth}p{0.47\linewidth}@{}}
a)~64 & b)~10 \\
c)~16 & d)~12 \\
\end{tabular}
\end{minipage}

\item \begin{minipage}[t]{\linewidth}
¿Cuál es el resultado de 15 \% 4?

\vspace{4pt}
\begin{tabular}{@{}p{0.47\linewidth}p{0.47\linewidth}@{}}
a)~3 & b)~1 \\
c)~3.75 & d)~4 \\
\end{tabular}
\end{minipage}

\item \begin{minipage}[t]{\linewidth}
¿Cuál es el resultado de 20 - 12 / 4 * 2?

\vspace{4pt}
\begin{tabular}{@{}p{0.47\linewidth}p{0.47\linewidth}@{}}
a)~4 & b)~18.5 \\
c)~1 & d)~14 \\
\end{tabular}
\end{minipage}

\item \begin{minipage}[t]{\linewidth}
Después de estas instrucciones, ¿cuánto vale precioFinal?

\vspace{2pt}{\small
\begin{verbatim}
CONSTANTE REAL IVA ← 0.15
REAL precioProducto ← 60.00
REAL precioFinal ← precioProducto + (precioProducto * IVA)
\end{verbatim}}


\vspace{4pt}
\begin{tabular}{@{}p{0.47\linewidth}p{0.47\linewidth}@{}}
a)~60.15 & b)~75 \\
c)~69 & d)~9 \\
\end{tabular}
\end{minipage}

\item \begin{minipage}[t]{\linewidth}
Después de estas instrucciones, ¿cuánto vale longitud?

\vspace{2pt}{\small
\begin{verbatim}
CONSTANTE REAL PI ← 3.1416
REAL r ← 5
REAL longitud ← 2 * PI * r
\end{verbatim}}


\vspace{4pt}
\begin{tabular}{@{}p{0.47\linewidth}p{0.47\linewidth}@{}}
a)~15.708 & b)~31.416 \\
c)~78.54 & d)~11.2832 \\
\end{tabular}
\end{minipage}

\item \begin{minipage}[t]{\linewidth}
Si x = 7 e y = 4, ¿cuál es el resultado de (x + y > 10) Y (x \% 2 == 0)?

\vspace{4pt}
\begin{tabular}{@{}p{0.47\linewidth}p{0.47\linewidth}@{}}
a)~1 & b)~verdadero \\
c)~11 & d)~falso \\
\end{tabular}
\end{minipage}

\item \begin{minipage}[t]{\linewidth}
Si A es verdadero y B es falso, ¿cuál de estas expresiones da verdadero?

\vspace{4pt}
\begin{tabular}{@{}p{0.47\linewidth}p{0.47\linewidth}@{}}
a)~NO A & b)~A O B \\
c)~NO (A O B) & d)~A Y B \\
\end{tabular}
\end{minipage}

\end{enumerate}
\end{document}
